\documentclass{beamer}
\usepackage[utf8]{inputenc}

%\usetheme{Copenhagen}
\usetheme{Antibes}
\hypersetup{pdfstartview={Fit}}

% Defining a new environment for continual enumerating throught frames
\makeatletter
\newenvironment{cenumerate}{
	\enumerate
	\setcounter{\@enumctr}{\csname saved@\@enumctr\endcsname}
}{
	\expandafter\xdef\csname saved@\@enumctr\endcsname{\the\value{\@enumctr}}
	\endenumerate
}
\newenvironment{cenumerate*}{
	\enumerate
}{
	\expandafter\xdef\csname saved@\@enumctr\endcsname{\the\value{\@enumctr}}
	\endenumerate
}
\makeatother

\title{Joint 3D object recognition and tracking}
\author[Orizio, Rizzini, Zucchelli]{Orizio Riccardo,\\Rizzini Mattia,\\Zucchelli Maurizio}
\date{25 july 2013}
\institute[UniBS]{University of Brescia}
\logo{\includegraphics[width=15mm]{logo.png}}

\begin{document}
	\begin{frame}
		\maketitle
	\end{frame}

	\section{Introduction}
	
	\begin{frame}
		\frametitle{\insertsection}
		The project goal is to realize a joint system of recognition and tracking
		of some features of a previously chosen object which is identified in a video.
		These features are called \textsc{Label}s.\\
		The sofware is aimed for execution on low performance devices such as the
		iPad, therefore the tracking is necessary due to the high computational
		cost of a pure recognition.\\
	\end{frame}

	\begin{frame}
		\frametitle{\insertsection\ - Overall structure}
		The software is designed for concurrent execution.\\
		There is an object of class \textsc{Manager} which, given every frame
		captured by a webcam or in a video, elaborates it in some way and retuns
		an object of class \textsc{Object} which contains the actual \textsc{Label}
		positions.\\
	\end{frame}

	\AtBeginSection[]
	{
		\begin{frame}
			\frametitle{Outline}
			\tableofcontents[currentsection]
		\end{frame}
	}
	\section{Manager}

	\begin{frame}
		\frametitle{\insertsection}
		The \textsc{Manager} handles two objects of class \textsc{Recognizer} and
		\textsc{Tracker}. The first object runs a \emph{thread}.\\
		Every $N$ frames elaborated, the \textsc{Manager} asks the
		\textsc{Recognizer} to perform a full recognition. In the meantime the
		\textsc{Tracker} keeps tracking the \textsc{Label}s starting from the last
		\textsc{Object} given by the \textsc{Recognizer}.\\
		When the \textsc{Recognizer} has finished it returns an \textsc{Object}
		which contains the positions of the \textsc{Label}s in the moment in which
		the recognition started. These positions are then \emph{actualized} by the
		\textsc{Tracker}.\\
	\end{frame}

	\section{Recognition}

	\begin{frame}
		\frametitle{\insertsection}
		To perform recognition, \emph{SIFT} descriptors are used since, in our tests,
		they seemed to perform better than the \emph{SURF} ones without any computational
		disadvantage.\\
	\end{frame}

	\subsection{Database}
	
	\begin{frame}
		\frametitle{\insertsection\ - \insertsubsection}
		The recognition bases itself on a \textsc{Database}.\\
		This \textsc{Database} supports serialization of the structures used for
		recognition. At the program start, a \textsc{Database} name is provided:
		if this name corresponds to a serialized \textsc{Database}, the
		\textsc{Database} is loaded; otherwise it is trained.\\
	\end{frame}

	\begin{frame}
		\frametitle{\insertsection\ - \insertsubsection}
		The \textsc{Database} training is performed through the analysis of a group
		of sample images depicting the object from various points of view.\\
		To every image is also associated a text file which contains the absolute
		positions of the various \textsc{Label}s in the referenced sample.\\
		Every sample image is loaded, from every image \emph{SIFT} features are
		extracted, descriptors computed and label positions loaded. Everything
		is then saved in three structures and serialized.\\
		At the end of the training or load phase, a \emph{FLANN} based matcher is
		trained based on all the samples descriptors.\\
	\end{frame}

	\subsection{Matching}

	\begin{frame}
		\frametitle{\insertsection\ - \insertsubsection}
		To perform a match on a frame, the following steps are executed:
		\begin{cenumerate*}
			\item The SIFT features and descriptors are extracted from the frame.
			\item The descriptors are matched using the \emph{FLANN} matcher with
				a \emph{knnMatch} method.
			\item The matches are searched to identify the sample which has got 
				the maximum number of matches. The only informations considered in
				the database are now those regarding this sample.
			\item The matches of that sample are extracted and filtered using the
				\emph{NNDR} criteria.
		\end{cenumerate*}
	\end{frame}

	\begin{frame}
		\frametitle{\insertsection\ - \insertsubsection}
		\begin{cenumerate}
			\item If the remaining matches are less than 20 (estimated practically)
				then the match ends with nothing found (empty \textsc{Object} is given back).
			\item Otherwise the keypoints related to the matches in the database and the
				frame are used to estimate an homography with \emph{RANSAC}.
			\item If it is valid (the inliers ratio over the total points used by RANSAC is
				at least 50\%), the homography is used to remap the positions of the labels in
				the frame. Those positions are added to an \textsc{Object}
				which is passed back to the \textsc{Manager}.
		\end{cenumerate}
	\end{frame}

	\section{Tracking}

	\begin{frame}
		\frametitle{\insertsection}
		The tracking of the object is divided in two principal parts:
		\begin{enumerate}
			\item The effective tracking of the \textsc{Object} between two
				successive frames.
			\item The \emph{actualization} of the \textsc{Object} received by the
				\textsc{Recognizer}.
		\end{enumerate}
	\end{frame}

	\begin{frame}
		\frametitle{\insertsection}
		To track an object, \emph{GFTT} (Good Features To Track) are used:
		these features are extremely fast to calculate and provide optimal results
		in tracking application.\\
		The algorithm works on the whole frame: the features are associated to the
		new frame using Lucas-Kanade's method for \emph{optical flow} estimation,
		then, for every \textsc{Label}, the nearest features to it are calculated
		using a \emph{FLANN} based matcher and their movement between the frames is
		mediated and added to the \textsc{Label}'s position.
	\end{frame}

	\subsection{Actualization}

	\begin{frame}
		\frametitle{\insertsection\ - \insertsubsection}
		The \emph{actualization} process updates the \textsc{Object} received from
		the \textsc{Recognizer} to its position in the current frame. This is
		necessary because of the long time required by it.\\
		To allow an easy and fast \emph{actualization}, the \textsc{Tracker}
		calculates and saves the features of the frame used for the recognition
		and starts using them internally for the tracking of the other frames.\\
		This way, when the \textsc{Object} is received, the only operation to be
		done is the tracking from the saved features to the current features.\\
	\end{frame}
	
	\subsection{Optimization}

	\begin{frame}
		\frametitle{\insertsection\ - \insertsubsection}
		The heaviest part of the tracking algorithm is the calculation of the \emph{optical
		flow}. To reduce the time required, the frame is decimated by a factor
		of $2$; since the complexity of the algorithm is more than linear, this
		reduces the computational cost to less than $25\%$ of the original time.\\
		Also, to furtherly reduce the complexity, the whole tracker has been
		adapted to work with decimated frames. This can be done without damaging
		the results because both the features and the tracking algorithm are robust
		to scaling.\\
	\end{frame}

	\section{Conclusions}

	\begin{frame}
		\frametitle{\insertsection}
		The software works very well in tracking simple shaped objects with
		lots of details (mostly because of the way SIFT works).\\
		The combination of a tracking and recognizing system achieved pretty
		well the goal of creating a lightwight software.
		It is to notice that to achieve this, some sacrifice has been made,
		particularly in the tracking part which has the purpouse to supply at
		the slowliness of the recognizer. The biggest consequence of this is
		the inaccuracy in tracking objects in fast motion.\\
		Since the requirements specified the tracking of slow motion, we
		are pretty good with that flaw.
	\end{frame}
\end{document}
